\section{硬件细节深讲：从一条指令到一次设备完成}

很多操作系统书厚，不是因为它们列了更多名词，而是因为它们会反复把同一个机制放到不同层次里解释。本节从硬件层展开，说明 OS 为什么必须关心指令执行、内存层次、设备协议和启动约束。

\subsection{一条 load 指令背后发生了什么}

用户程序写下一句 \code{x = *p}，看起来只是读取内存。硬件实际要经过多个检查和缓存层次。CPU 先得到虚拟地址，TLB 查询是否有缓存的翻译；若 TLB 未命中，硬件页表遍历器读取多级页表；页表项必须 present，并允许当前特权级读取；然后物理地址进入 cache 层次，L1 未命中查 L2，L2 未命中查 L3，最后才可能访问内存控制器。

\begin{lstlisting}
load virtual_address:
  translation = tlb_lookup(virtual_address)
  if translation.miss:
    pte = page_table_walk(current.cr3, virtual_address)
    if not pte.present or not pte.readable:
      raise_page_fault()
    tlb_insert(virtual_address, pte)
  physical = pte.physical_page + page_offset(virtual_address)
  line = cache_lookup(physical)
  if line.miss:
    line = memory_read(cache_line_address(physical))
  return extract_word(line, physical)
\end{lstlisting}

这条路径解释了几个 OS 事实。第一，页表权限是硬件执行的，不是内核每次手工检查。第二，TLB 让地址翻译变快，但也让页表修改必须配合失效。第三，cache line 让相邻数据访问更快，也让伪共享和锁竞争变得昂贵。第四，缺页异常不是函数调用，而是 CPU 在执行普通指令时发现无法继续，转入内核处理。

\subsection{一次 store 为什么需要内存屏障}

普通程序通常假设代码顺序就是别人看到的顺序。多核和设备环境中，这个假设不总成立。CPU 可以把 store 暂存在 store buffer，编译器可以重排不相关内存访问，设备可能通过总线异步观察内存。内核因此使用 memory barrier 表达顺序约束。

\begin{lstlisting}
publish_object(obj):
  obj.field1 = value1
  obj.field2 = value2
  store_release(global_ptr, obj)

read_object():
  obj = load_acquire(global_ptr)
  if obj != null:
    use(obj.field1, obj.field2)
\end{lstlisting}

release/acquire 的意思是：发布者在指针可见前完成对象初始化；读取者看到指针后，也能看到发布者之前写入的字段。锁、RCU、无锁队列和驱动描述符都依赖类似规则。没有屏障，bug 可能只在特定 CPU、特定优化级别、特定负载下出现。

\subsection{NUMA：内存不是均匀的}

多路服务器上，CPU 插槽和内存控制器形成 NUMA node。本地内存访问比远端内存低延迟、高带宽。调度器、页分配器和内存回收因此不能只把所有 CPU 和内存看成一个平面池。

\begin{longtable}{p{0.22\linewidth}p{0.34\linewidth}p{0.32\linewidth}}
\toprule
现象 & 原因 & OS 对策 \\
\midrule
线程迁移后变慢 & cache 和 TLB 热状态丢失 & CPU affinity、迁移成本估计 \\
访问远端内存慢 & 内存挂在另一个 NUMA node & first-touch、NUMA-aware allocation \\
局部内存耗尽 & 某 node 压力更高 & 本地回收、跨 node fallback \\
锁竞争放大 & cache line 在插槽间移动 & per-CPU 数据、分片锁、RCU \\
\bottomrule
\end{longtable}

这也是为什么“线程数等于核心数”不是万能规则。若线程频繁共享同一批 cache line，即使核心很多，也可能被 cache coherence 拖慢。

\subsection{MMIO、PIO 和 DMA 的差别}

设备 I/O 有三种典型方式。PIO 由 CPU 直接搬运数据，适合简单设备或早期启动。MMIO 让 CPU 通过地址读写设备寄存器，适合控制设备状态。DMA 让设备直接访问内存，适合大吞吐数据传输。

\begin{longtable}{p{0.18\linewidth}p{0.32\linewidth}p{0.38\linewidth}}
\toprule
方式 & 数据路径 & OS 难点 \\
\midrule
PIO & CPU 读写数据端口或寄存器 & CPU 被搬运占用，吞吐低 \\
MMIO & CPU 读写设备寄存器 & 访问顺序、副作用、超时、错误状态 \\
DMA & 设备直接读写内存 & buffer pin、IOMMU、cache 一致性、completion \\
\bottomrule
\end{longtable}

一个块设备写请求可能同时使用三者：驱动用 MMIO 配置队列，用 DMA 传输数据，用中断报告 completion。OS 教材如果只讲“设备读写”，就会漏掉驱动里最容易出错的生命周期和顺序问题。

\subsection{启动时为什么不能使用完整内核能力}

启动早期没有堆、没有调度、没有完整页表、没有普通文件系统，也不一定有中断。很多初学者会问：为什么不一开始就调用统一日志、统一分配器、统一驱动框架？原因是这些框架本身依赖更基础的能力。

\begin{lstlisting}
early code can use:
  fixed stack
  linker symbols
  bootloader memory map
  early console
  bump allocator

early code cannot assume:
  scheduler exists
  page allocator is ready
  interrupts are enabled
  file system is mounted
  all CPUs are online
\end{lstlisting}

启动代码的复杂性不是历史负担，而是依赖关系决定的。一个成熟内核会把初始化分层：先建立能输出错误的 console，再建立能描述内存的 boot memory，再建立页表和分配器，最后建立调度器、设备和用户态。

\subsection{硬件基础检查题}

\begin{enumerate}
\tightlist
\item 为什么修改页表后还要处理 TLB，而不是只改内存中的页表项？
\item 为什么驱动写完描述符后要有 memory barrier 再通知设备？
\item 为什么中断处理函数不能随意睡眠？
\item 为什么 DMA buffer 在 completion 前不能释放？
\item 为什么多核调度不能只用一个全局 runqueue？
\item 为什么启动早期不能调用普通 \code{malloc} 或普通文件 I/O？
\end{enumerate}

能回答这些问题，后面的虚拟内存、调度、驱动和文件系统才有硬件根基。

\subsection{时间尺度：为什么 OS 总在权衡}

操作系统的大量设计都是在不同时间尺度之间权衡。寄存器访问可以看作 1 个 CPU 周期量级，L1 cache 是几个周期，内存可能是上百周期，NVMe I/O 可能是微秒到毫秒，网络和机械磁盘可能更慢。调度器一次上下文切换、一次 TLB shootdown、一次 fsync 都不是“抽象操作”，它们背后有硬件成本。

\begin{longtable}{p{0.22\linewidth}p{0.28\linewidth}p{0.38\linewidth}}
\toprule
操作 & 典型尺度 & OS 设计影响 \\
\midrule
寄存器访问 & CPU 周期 & 上下文切换要保存恢复寄存器 \\
L1/L2 cache 命中 & 几到几十周期 & 数据局部性、per-CPU 数据结构重要 \\
主内存访问 & 上百周期量级 & page fault 和 cache miss 成本明显 \\
TLB shootdown & 跨 CPU 中断和等待 & 频繁改页表会拖慢多核程序 \\
NVMe I/O & 微秒到毫秒 & 需要队列、异步、批量和 completion \\
fsync/flush & 可能毫秒级或更高 & 可靠提交不能放在热路径随意调用 \\
\bottomrule
\end{longtable}

这些数字不是让读者背精确值，而是建立数量级直觉。一个算法少一次内存访问可能值得优化；一个功能如果每个请求都 fsync，就可能把吞吐限制在存储设备 flush 延迟上；一个锁如果让 cache line 在 CPU 间来回迁移，会比临界区本身更贵。

\subsection{一个完整硬件路径：网卡收到一个包}

以网卡接收数据包为例，可以把硬件和 OS 串起来。

\begin{lstlisting}
1. driver allocates RX buffers
2. driver maps buffers for DMA
3. driver fills RX descriptor ring
4. NIC receives packet from wire
5. NIC DMA-writes packet into memory
6. NIC writes completion/status
7. NIC raises interrupt or waits for interrupt moderation
8. interrupt handler schedules poll
9. poll reads descriptors, invalidates cache if needed
10. network stack consumes packet
11. socket receive queue wakes user process
\end{lstlisting}

这条路径里任何一步都可能出错：buffer 太少会丢包；DMA 映射错误会写坏内存；中断合并过大增加延迟；poll budget 太大影响调度；用户进程不读 socket 会导致接收队列积压并产生背压。教材要讲厚，就必须把这样的路径完整展开，而不是只说“网卡通过中断通知内核”。

\subsection{从单片机循环到通用操作系统}

理解操作系统最好不要从“Linux 有哪些子系统”开始，而要从最小机器开始。一个简单单片机程序通常只有一个主循环和若干中断处理函数。它没有用户态和内核态之分，没有页表，没有进程，也没有文件系统。程序直接读写外设寄存器，所有代码处在同一信任边界内。

\begin{lstlisting}
bare_metal_main():
  init_clock()
  init_gpio()
  init_uart()
  while true:
    poll_button()
    update_led()
    handle_uart_bytes()
\end{lstlisting}

这种模型能工作，是因为问题规模小：任务数量少，开发者知道所有代码，内存布局固定，失败影响有限。一旦系统要同时运行多个互不信任程序，或要让一个程序卡住时不拖死全部功能，就需要更强的机制。

\begin{longtable}{p{0.2\linewidth}p{0.36\linewidth}p{0.32\linewidth}}
\toprule
阶段 & 新需求 & 产生的 OS 机制 \\
\midrule
单循环裸机 & 周期性处理少量设备 & 主循环、轮询、少量中断 \\
中断驱动裸机 & 外设事件不能丢 & 中断向量、临界区、环形缓冲区 \\
协作式多任务 & 多个任务共享 CPU & 任务控制块、显式 yield、简单栈切换 \\
抢占式内核 & 任务不能长期霸占 CPU & 定时器中断、调度器、保存恢复上下文 \\
多程序系统 & 程序互不信任 & 特权级、系统调用、页表、权限检查 \\
通用系统 & 设备和文件很多 & 驱动模型、VFS、块层、网络栈、日志恢复 \\
\bottomrule
\end{longtable}

这条演进线解释了为什么 OS 特性不是凭空增加的。调度器来自“不能信任任务主动让出 CPU”；系统调用来自“用户程序不能直接改内核状态”；虚拟内存来自“程序不能互相读写”；文件系统来自“存储设备不能只暴露原始扇区”；网络栈来自“设备字节流必须变成可复用通信接口”。

\subsection{协作式任务为什么不够}

协作式任务把每个任务写成能主动让出的函数。它比单循环清晰，但它依赖每个任务都守规矩。一个任务忘记调用 \code{yield}，或者在循环中等待硬件状态，就会让系统停止响应。

\begin{lstlisting}
task_a():
  while true:
    do_a_small_step()
    yield()

task_b_bug():
  while device_not_ready():
    ;   // forgot to yield, all other tasks starve
\end{lstlisting}

抢占式调度用定时器中断打破这个假设。即使当前任务不合作，硬件也会在时间片到达时进入内核。内核保存当前寄存器，选择另一个 runnable 任务，再恢复其寄存器。这就是“硬件定时器 + trap frame + scheduler”共同形成的控制权回收机制。

\subsection{trap frame：异常入口的最小证据}

CPU 进入内核时，必须留下足够信息，让内核知道从哪里来、为什么来、如何返回。这个保存下来的状态通常称为 trap frame 或 interrupt frame。不同架构字段不同，但语义类似。

\begin{longtable}{p{0.22\linewidth}p{0.34\linewidth}p{0.32\linewidth}}
\toprule
字段 & 含义 & 用途 \\
\midrule
RIP & 被打断或出错的指令位置 & 返回、信号、调试栈 \\
RSP & 进入异常前的栈指针 & 恢复用户栈或诊断坏栈 \\
RFLAGS & 中断位、条件码、特权级信息 & 判断返回权限和中断状态 \\
通用寄存器 & 参数和临时值 & 系统调用参数、恢复用户执行 \\
错误码 & page fault、保护异常等附加原因 & 区分缺页、权限和非法访问 \\
\bottomrule
\end{longtable}

trap frame 解释了系统调用和异常为何能返回到原程序。它也解释了为什么内核 panic 时要打印寄存器和指令地址：这些字段是定位错误路径的最低成本证据。

\begin{lstlisting}
trap_entry(vector, hardware_frame):
  switch_to_kernel_stack_if_from_user()
  save_general_registers()
  tf = current.trap_frame
  dispatch_trap(vector, tf)
  restore_general_registers()
  return_from_trap(tf)
\end{lstlisting}

入口代码必须非常保守。它运行在栈切换、页表边界和中断状态都敏感的位置，不能调用可能睡眠的复杂函数。很多内核把低级入口代码写成汇编或极薄的 C 包装，就是为了精确控制保存顺序和 ABI。

\subsection{中断控制器、优先级和 IPI}

设备不是直接“喊 CPU”。多数机器有中断控制器负责接收设备中断、屏蔽中断、设置优先级并把中断送到某个 CPU。多核系统还需要 CPU 之间互相发送 IPI，用于重调度、TLB shootdown、停止其他 CPU 或调试。

\begin{longtable}{p{0.24\linewidth}p{0.34\linewidth}p{0.3\linewidth}}
\toprule
机制 & 用途 & OS 风险 \\
\midrule
interrupt mask & 临时屏蔽某类中断 & 屏蔽过久造成延迟或丢事件 \\
interrupt priority & 高优先级事件先处理 & 优先级反转或低优先级饥饿 \\
interrupt affinity & 把设备中断送到指定 CPU & 热点 CPU 过载或缓存局部性变差 \\
MSI/MSI-X & 设备用内存写触发中断消息 & 向量分配、隔离和设备错误处理 \\
IPI & CPU 间发送控制事件 & 频繁 IPI 会破坏多核伸缩性 \\
\bottomrule
\end{longtable}

TLB shootdown 是 IPI 的典型用途。一个 CPU 修改了某进程页表，其他正在运行该进程的 CPU 可能仍有旧 TLB 项。修改方必须通知这些 CPU 失效旧翻译，并等待确认后才能释放旧物理页。这个过程昂贵，所以现代内核会尽量批量处理 unmap 和权限变化。

\subsection{内存模型：编译器顺序、CPU 顺序和设备顺序}

很多并发错误来自把三种顺序混在一起。源代码顺序是程序员写下的顺序；编译器可能在不改变单线程语义的前提下重排；CPU 可能因为 store buffer、乱序执行和缓存协议让其他 CPU 以不同时间观察到写入；设备通过 DMA 和 MMIO 观察到的顺序还要额外受总线和设备协议影响。

\begin{longtable}{p{0.2\linewidth}p{0.38\linewidth}p{0.3\linewidth}}
\toprule
顺序层次 & 例子 & 常用约束 \\
\midrule
编译器顺序 & 普通变量访问被优化或重排 & volatile 语义、编译器屏障、原子操作 \\
CPU 顺序 & store 先进入缓冲区，稍后对其他核可见 & acquire/release、full fence \\
设备顺序 & 描述符写入和 doorbell MMIO 顺序 & DMA barrier、MMIO write barrier \\
\bottomrule
\end{longtable}

驱动发布 DMA 描述符时，必须先让内存中的描述符对设备可见，再写 MMIO doorbell。锁释放时，必须先让临界区写入对其他 CPU 可见，再让锁状态变成未持有。RCU 发布指针时，必须先初始化对象，再发布地址。这些看似不同的场景，本质都是“先写数据，再发布可见性”的顺序要求。

\subsection{硬件抽象的边界}

OS 会抽象硬件，但不能把硬件成本抹掉。页表让每个进程看到独立地址空间，但 TLB shootdown 仍然昂贵；block layer 把设备包装成请求队列，但 flush 仍然慢；socket 把网卡包装成字节流，但拥塞和丢包仍然存在；调度器把 CPU 时间切片给线程，但 cache/TLB 热状态仍然影响性能。

\begin{rubricbox}
学习硬件基础的验收标准：看到一个 OS 机制时，能回答它依赖哪个硬件入口、改变哪些硬件状态、失败时硬件给出什么证据、性能瓶颈来自哪个硬件层次。
\end{rubricbox}

\subsection{本章小结}

\begin{itemize}
\tightlist
\item CPU 特权级和页表权限是 OS 隔离的硬件根。
\item cache、TLB、NUMA 和内存屏障解释了内核为何重视局部性和顺序。
\item MMIO、DMA、中断和 IOMMU 构成设备驱动的基本硬件接口。
\item 启动阶段受依赖限制，不能使用尚未初始化的内核服务。
\item 所有高级 OS 机制都能追溯到硬件状态、硬件事件和硬件成本。
\end{itemize}
